% GENERATED FILE. Edit canonical lessons, then run npm run generate:books.
% canonical-source-hash: fnv1a64:894e2b3725816686
% canonical-lessons: MR-C101-undir, MR-C101-simha, MR-C101-vagh, MR-C101-hatti, MR-C101-makad

\chapter{Mouse, Lion, Tiger, Elephant, Monkey}
\label{ch:mr-mouse-lion-tiger-elephant-monkey}
\hlchaptermodality{101}

\begin{chapteropening}
\textbf{By the end of this chapter:} \emph{I can say a mouse, a lion, a tiger, an elephant and a monkey.}

Complete the last lesson of chapter 101: I can say a mouse, a lion, a tiger, an elephant and a monkey.
\end{chapteropening}

\section[\texorpdfstring{\mr{उंदीर} (undīr)}{उंदीर (undīr)}]{\mr{उंदीर} (undīr) — a mouse}

\label{lesson:MR-C101-undir}



\begin{quote}
Before the new one: say the Marathi for a sheep, then the Marathi for a bird.
\end{quote}

\subsection*{You'll want to know: \mr{उंदीर}}
\textbf{\mr{उंदीर}} — \emph{undīr} — \textquotedblleft{}a mouse\textquotedblright{}.

\begin{cousinweb}
\textbf{\mr{उंदीर}} (\emph{undīr}) is a tadbhava of Sanskrit \emph{undura}, \textquotedblleft{}rat, mouse\textquotedblright{} (CDIAL 2095). CDIAL notes the view that the Sanskrit word comes from Austroasiatic.
\end{cousinweb}

\begin{grammarlens}[title={What you've built}]
One more word for this chapter.
\end{grammarlens}

\subsection*{Your turn}
\begin{itemize}
\raggedright
  \item \emph{Say it:} \emph{undīr}
  \item \emph{Say it:} \emph{undīr}, once more
  \item \emph{Say it:} say \emph{undīr}
  \item \emph{Recall:} say the Marathi for a sheep, then the Marathi for a bird, then say \emph{undīr} again
\end{itemize}

\begin{tcolorbox}[breakable,colback=teal!4,colframe=teal!35!black,title={Before you move on}]
What does \mr{उंदीर} mean? (\textquotedblleft{}A mouse\textquotedblright{}.) Say it once more.
\end{tcolorbox}
\section[\texorpdfstring{\mr{सिंह} (siṁha)}{सिंह (siṁha)}]{\mr{सिंह} (siṁha) — a lion}

\label{lesson:MR-C101-simha}



\begin{quote}
Before the new one: say the Marathi for a bird, then the Marathi for a mouse.
\end{quote}

\subsection*{You'll want to know: \mr{सिंह}}
\textbf{\mr{सिंह}} — \emph{siṁha} — \textquotedblleft{}a lion\textquotedblright{}.

\begin{cousinweb}
\textbf{\mr{सिंह}} (\emph{siṁha}) is Sanskrit \emph{siṁha}, \textquotedblleft{}lion\textquotedblright{}, taken over whole as a learned loan, a tatsama (CDIAL 13384).
\end{cousinweb}

\begin{grammarlens}[title={What you've built}]
One more word for this chapter.
\end{grammarlens}

\subsection*{Your turn}
\begin{itemize}
\raggedright
  \item \emph{Say it:} \emph{siṁha}
  \item \emph{Say it:} \emph{siṁha}, once more
  \item \emph{Say it:} say \emph{siṁha}
  \item \emph{Recall:} say the Marathi for a bird, then the Marathi for a mouse, then say \emph{siṁha} again
\end{itemize}

\begin{tcolorbox}[breakable,colback=teal!4,colframe=teal!35!black,title={Before you move on}]
What does \mr{सिंह} mean? (\textquotedblleft{}A lion\textquotedblright{}.) Say it once more.
\end{tcolorbox}
\section[\texorpdfstring{\mr{वाघ} (vāgh)}{वाघ (vāgh)}]{\mr{वाघ} (vāgh) — a tiger}

\label{lesson:MR-C101-vagh}



\begin{quote}
Before the new one: say the Marathi for a mouse, then the Marathi for a lion.
\end{quote}

\subsection*{You'll want to know: \mr{वाघ}}
\textbf{\mr{वाघ}} — \emph{vāgh} — \textquotedblleft{}a tiger\textquotedblright{}.

\begin{cousinweb}
\textbf{\mr{वाघ}} (\emph{vāgh}) is a tadbhava, a word inherited from Sanskrit and reshaped by sound change. It comes from \emph{vyāghra}, \textquotedblleft{}tiger\textquotedblright{} (CDIAL 12193).
\end{cousinweb}

\begin{grammarlens}[title={What you've built}]
One more word for this chapter.
\end{grammarlens}

\subsection*{Your turn}
\begin{itemize}
\raggedright
  \item \emph{Say it:} \emph{vāgh}
  \item \emph{Say it:} \emph{vāgh}, once more
  \item \emph{Say it:} say \emph{vāgh}
  \item \emph{Recall:} say the Marathi for a mouse, then the Marathi for a lion, then say \emph{vāgh} again
\end{itemize}

\begin{tcolorbox}[breakable,colback=teal!4,colframe=teal!35!black,title={Before you move on}]
What does \mr{वाघ} mean? (\textquotedblleft{}A tiger\textquotedblright{}.) Say it once more.
\end{tcolorbox}
\section[\texorpdfstring{\mr{हत्ती} (hattī)}{हत्ती (hattī)}]{\mr{हत्ती} (hattī) — an elephant}

\label{lesson:MR-C101-hatti}



\begin{quote}
Before the new one: say the Marathi for a lion, then the Marathi for a tiger.
\end{quote}

\subsection*{You'll want to know: \mr{हत्ती}}
\textbf{\mr{हत्ती}} — \emph{hattī} — \textquotedblleft{}an elephant\textquotedblright{}.

\begin{cousinweb}
\textbf{\mr{हत्ती}} (\emph{hattī}) is a tadbhava of Sanskrit \emph{hastin}, \textquotedblleft{}having hands\textquotedblright{}, which meant \textquotedblleft{}elephant\textquotedblright{} from the Rigveda on (CDIAL 14039). CDIAL suggests it renders \emph{karin}, \textquotedblleft{}elephant\textquotedblright{}, a word perhaps of non-Aryan origin, linked by folk etymology with \emph{kara}, \textquotedblleft{}hand\textquotedblright{}.
\end{cousinweb}

\begin{grammarlens}[title={What you've built}]
One more word for this chapter.
\end{grammarlens}

\subsection*{Your turn}
\begin{itemize}
\raggedright
  \item \emph{Say it:} \emph{hattī}
  \item \emph{Say it:} \emph{hattī}, once more
  \item \emph{Say it:} say \emph{hattī}
  \item \emph{Recall:} say the Marathi for a lion, then the Marathi for a tiger, then say \emph{hattī} again
\end{itemize}

\begin{tcolorbox}[breakable,colback=teal!4,colframe=teal!35!black,title={Before you move on}]
What does \mr{हत्ती} mean? (\textquotedblleft{}An elephant\textquotedblright{}.) Say it once more.
\end{tcolorbox}
\section[\texorpdfstring{\mr{माकड} (mākaḍ)}{माकड (mākaḍ)}]{\mr{माकड} (mākaḍ) — a monkey}

\label{lesson:MR-C101-makad}



\begin{quote}
Before the new one: say the Marathi for a tiger, then the Marathi for an elephant.
\end{quote}

\subsection*{You'll want to know: \mr{माकड}}
\textbf{\mr{माकड}} — \emph{mākaḍ} — \textquotedblleft{}a monkey\textquotedblright{}.

\begin{cousinweb}
\textbf{\mr{माकड}} (\emph{mākaḍ}) is a tadbhava of Sanskrit \emph{markaṭa}, \textquotedblleft{}monkey\textquotedblright{} (CDIAL 9882). CDIAL suggests it may extend a word \emph{marka} that may come from Dravidian, comparing Kannada \emph{maṅga}, \textquotedblleft{}monkey\textquotedblright{}.
\end{cousinweb}

\begin{grammarlens}[title={What you've built}]
One more word for this chapter.
\end{grammarlens}

\subsection*{Your turn}
\begin{itemize}
\raggedright
  \item \emph{Say it:} \emph{mākaḍ}
  \item \emph{Say it:} \emph{mākaḍ}, once more
  \item \emph{Say it:} say \emph{mākaḍ}
  \item \emph{Recall:} say the Marathi for a tiger, then the Marathi for an elephant, then say \emph{mākaḍ} again
\end{itemize}

\begin{tcolorbox}[breakable,colback=teal!4,colframe=teal!35!black,title={Before you move on}]
What does \mr{माकड} mean? (\textquotedblleft{}A monkey\textquotedblright{}.) Say it once more.
\end{tcolorbox}
